\section{Introduction}
\label{sec:intro}


%-------------------------------------------------------------------------
\subsection{Task Description}

The machine learning competition focused on top-k similar images retrieval, a relevant computer vision task with many applications across different domains.
Participants were challenged to develop models capable of identifying and ranking the 10 most similar images to a given query from a large database.
The core objective was to retrieve the most similar synthetic images from a gallery when given natural query images of celebrities.
The training set contained approximately 5000 images organized into subfolders by celebrity identity, i.e. with each subfolder containing both natural and synthetic images of the same individual.
The test set, on the other hand, presented a more constrained scenario: query set consisted of approximately 3000 real images of celebrities, while the gallery contained only synthetic ones. 
This set up (which was unknown before the competition day) tested models' ability to connect authentic photographs and AI-generated synthetic representations of the same celebrities (and in most cases, the same photo too). 
For each query image, models were required to retrieve 10 gallery images ranked by similarity, with the most similar image ranked first. 
The challenge, in fact, laid not only in identifying the correct celebrity but also in properly ordering the retrieved results based on visual similarity confidence. 

\subsection{Overview of approaches}

Our team explored multiple architectural paradigms to address this cross-domain retrieval challenge, employing both traditional computer vision approaches and modern foundantion models.
\begin{itemize}
\item {\bf Traditional CNN Architectures}: starting from traditional convolutional neural networks (including EfficientNet and ResNet), we implemented fine-tuning strategies.
These architectures were adapted from their original image classification objectives to learn similarity embeddings suitable for retrieval tasks. 
\item {\bf Transformer-based Models}: given that the details of the task were unknown until the day of the competition, the team expanded the investigation to include transformer-based architectures. 
We implemented DINO, a self supervised vision transformer that learns robust visual representations through knowledge distillation. 
Additionally, we employed CLIP (Contrastive Language-Image Pre-training), specifically the RN50x64 variant (known for the strong performance in visual tasks), which creates joint vision-language embeddings through contrastive learning on large-scale image-text pairs. 
\end{itemize}
Our approach evolved from conventional supervised fine-tuning towards more sophisticated methods, recognizing the limitations of traditional CNNS for cross-domain tasks, ultimately focusing on models that were pre-trained on diverse data that could capture semantic similarity beyond pixel-level correspondance. 

\subsection{Summary of results }
The competition revealed significant differences in performance across architectural choices. 
Despite extensive fine-tuning efforts on the provided training data, both EfficientNet and ResNet architectures demonstrated poor performance across all evaluation metrics. 
These models struggled significantly with the natural-to-synthetic domain transfer, suggesting that traditional convolutional features learned through supervised fine-tuning were insufficient for bridging the substantial visual gap between authentic photographs and AI-generated celebrity images.
During the competition, CLIP RN50x64 emerged as the decisive winner, consistently and substantially outperforming all other implemented approaches across all evaluation metrics. 
Indeed, CLIP RN50x60 achieved a score of 734.693/1000, summing up scores in Top-1 Accuracy (600 points), Top-5 Accuracy (300 points), and Top-10 Accuracy (100 points).
This superior performance can be attributed to CLIP's unique contrastive pre-training methodology, which learns to associate visual and textual representations across massive datasets, resulting in more semantically meaningful embeddings that generalize effectively across visual domains.


These results highlight a fundamental shift in computer vision capabilities, where foundation models (i.e. an AI neural network — trained on mountains of raw data, generally with unsupervised learning — that can be adapted to accomplish a broad range of tasks) trained on diverse, large-scale data with sophisticated objectives substantially outperform traditional architectures, even when the latter are specifically fine-tuned for the target domain. 
The findings emphasize the critical importance of pre-training methodology and semantic understanding for challenging cross-domain retrieval tasks in modern computer vision applications.

Code and model configurations presented in this paper are available at: \url{https://github.com/brunelliMichele/py.tatine}

